% Section 3.3, from lectures/L03/README.md: the objectives, the evaluation questions, and the next
% lecture.

\needspace{10\baselineskip}
\section{Review}
\label{c3:sec:review}

\subsection*{What you should be able to do now}
After this chapter, you should:
\begin{itemize}
  \item Know what a neural network is and be able to explain concepts such as layers, nodes, and
    weights.
  \item Be able to train small neural networks by hand.
  \item Have created an interface for a dense layer and a stub implementation of it.
\end{itemize}

\needspace{6\baselineskip}
\subsection*{Questions to test yourself}
You should be able to explain:
\begin{enumerate}
  \item Can you, in your own words, explain what happens in a neural network during feedforward and
    backpropagation?
  \item What's the purpose of an activation function, and what happens if you don't use one?
  \item What happens to training if the learning rate is set too high or too low?
  \item Why is it beneficial to define an interface for the dense layer rather than implementing a
    class directly?
  \item Your stub implements \code{initParams()} as an empty method. Why does the interface declare
    a method the stub has nothing to do for, and what would the network in
    Chapter~\ref{c4:ch:network} lose without it?
\end{enumerate}

\needspace{6\baselineskip}
\subsection*{Looking ahead}
Chapter~\ref{c4:ch:network} continues the exercise:
\begin{itemize}
  \item A neural network interface, and a concrete \code{Shallow} class that uses the stub you just
    built.
\end{itemize}
